\documentclass[12pt, a4paper]{article}

\usepackage[utf8]{inputenc}
\usepackage[T1]{fontenc}
\usepackage[italian]{babel}

\usepackage{geometry}
\geometry{a4paper, margin=2.5cm}

\usepackage{enumitem}

\title{\textbf{Abstract di Progetto: Shipping-on-the-Air}\\ \large (B2B)}
\author{}
\date{}

\begin{document}

\maketitle

\section*{Il Problema e l'Opportunità di Business }
Un'azienda cliente, attualmente limitata dalla logistica tradizionale su gomma (soggetta a ritardi, traffico urbano e alti costi operativi), ha la necessità strategica di scalare le proprie operazioni e ampliare la portata dei suoi servizi. L'opportunità di business consiste nel fornire a questa azienda un vantaggio competitivo esclusivo, permettendole di superare i limiti infrastrutturali terrestri sfruttando lo spazio aereo cittadino per l'ultimo miglio.

\section*{La Soluzione Proposta (Project Goal)}
Il progetto, commissionato specificamente per supportare l'espansione del cliente, mira a definire e gestire il rilascio di ``Shipping-on-the-Air'': un sistema integrato di consegne automatizzate tramite una flotta di droni. L'obiettivo è permettere all'azienda cliente di offrire consegne ultra-rapide (entro 30 minuti), riducendo i costi della propria flotta e offrendo un livello di servizio innovativo ai propri utenti finali.

\section*{Scope e Sottosistemi}
Per garantire la perfetta integrazione con il modello di business dell'azienda cliente, il progetto prevede la gestione di sei macro-funzionalità:
\begin{itemize}[label=\textbullet]
    \item \textbf{Gestione Ordini (Delivery \& Order):} Ricezione della richiesta e orchestrazione del ciclo di vita della consegna.
    \item \textbf{Logistica e Flotta (Drone Fleet):} Verifica dei vincoli tecnici pre-volo (autonomia batteria, rispetto della capacità massima di peso/payload).
    \item \textbf{Navigazione (Routing \& GPS):} Calcolo e assegnazione delle rotte di volo ottimali.
    \item \textbf{Tracciamento Live (Tracking):} Sistema di monitoraggio in tempo reale.
    \item \textbf{Notifiche (Notification):} Invio di avvisi automatici durante i cambi di stato della consegna (decollo, arrivo, ecc.).
    \item \textbf{Pagamenti (Payment):} Gestione transazioni ed eventuali rimborsi automatici in caso di fallimento.
\end{itemize}

\section*{Impostazione di Project Management}
Essendo l'obiettivo la simulazione della gestione progettuale, il focus si concentrerà su:
\begin{itemize}[label=\textbullet]
    \item \textbf{Gestione degli Stakeholder:} Forte coinvolgimento del cliente aziendale (Sponsor) per assicurarsi che i requisiti soddisfino le sue specifiche Conditions of Satisfaction.
    \item \textbf{Scelta del PMLC:} Adozione di un modello di ciclo di vita Iterativo o Adattivo per gestire l'incertezza tecnologica e rispondere prontamente ai feedback del cliente.
    \item \textbf{Gestione dei Rischi (Risk Management):} Mitigazione proattiva di rischi esterni e tecnologici (es.\ perdita segnale GPS, variazioni meteo, aggiornamenti normativi).
\end{itemize}

\end{document}